% ---- pg-428 （书页 416） ----

极大似然解. 解释这些因子.

\noindent 9.19\quad 一家公司正试图对其销售员工的质量作评估, 并且正寻找一种考察或一系列测试, 以期可以揭示是否有创造良好销售额的潜能. 该公司已挑选了 50 个销售人员的随机样本, 还已对每一个人就 3 项表现作了评估: 销售增长、销售利润和新客户销售额. 这些测度量已被变为同一尺度, 其中 100 表示“平均”业绩. 50 个人中的每一个接受 4 项测试, 分别测量创造力、机械推理、抽象推理和数学能力. \hyperref[tab:9.12]{表 9.12} 中列出在 $p=7$ 个变量上的 $n=50$ 次观察值.

(a) 假设对标准化变量 $Z_i=(X_i-\mu_i)/\sqrt{\sigma_{ii}}$ ($i=1,2,\cdots,7$) 有正交因子模型. 求 $m=2$ 和 $m=3$ 个公共因子的主成分解或极大似然解.

(b) 由(a)的解, 求 $m=2$ 和 $m=3$ 的旋转载荷. 比较这两组旋转载荷. 解释 $m=2$ 和 $m=3$ 的因子解.

(c) 列出 $m=2$ 和 $m=3$ 的共性方差、特殊方差和 $\hat{\mathbf{L}}\hat{\mathbf{L}}'+\hat{\boldsymbol{\Psi}}$ 的估计值. 比较这些结果. 此时, 你更愿意选择 $m$ 等于什么值? 为什么?

(d) 对 $m=2$ 和 $m=3$ 分别在 $\alpha=0.01$ 的水平下作检验 $H_0:\boldsymbol{\Sigma}=\mathbf{L}\mathbf{L}'+\boldsymbol{\Psi}$ 对 $H_1:\boldsymbol{\Sigma}\neq\mathbf{L}\mathbf{L}'+\boldsymbol{\Psi}$. 由这些结果及(b)和(c)中结果, 选择 $m$ 为多少最好?

(e) 设随机选取一个新的销售人员, 得到测验分数 $\mathbf{x}'=[x_1,x_2,\cdots,x_7]=[110,98,105,15,18,20,35]$. 用加权最小二乘法和回归方法, 计算这个销售人员的因子得分.
注意: $\mathbf{x}$ 的成分必须用由原始数据计算的样本均值和方差来标准化.

\noindent 9.20\quad 用\hyperref[tab:1.5]{表 1.5} 给出的空气污染变量 $X_1,X_2,X_5$ 和 $X_6$, 生成样本协方差矩阵.
(a) 求 $m=1$ 和 $m=2$ 的因子模型的主成分解.
(b) 求 $m=1$ 和 $m=2$ 时 $\mathbf{L}$ 和 $\boldsymbol{\Psi}$ 的极大似然估计.
(c) 比较由主成分和极大似然方法解得到的因子分解.

\noindent 9.21\quad 对练习 9.20 中 $m=2$ 的两个解都作最大方差旋转. 解释你得到的结果. 主成分分解和极大似然解相互一致吗?

\noindent 9.22\quad 参看练习 9.20.
(a) 用(i)\hyperref[eq:9-50]{式(9-50)}中的加权最小二乘和(ii)\hyperref[eq:9-58]{式(9-58)}的回归方法, 计算 $m=2$ 的极大似然估计的因子得分.
(b) 用\hyperref[eq:9-51]{式(9-51)}求主成分分解的因子得分.
(c) 比较此三组因子得分.

\noindent 9.23\quad 从样本相关系数矩阵着手, 重复练习 9.20. 解释 $m=1$ 和 $m=2$ 解的因子. 若 $\mathbf{R}$ 而非 $\mathbf{S}$ 因子化, 结果会有不同吗? 解释一下.

\noindent 9.24\quad 对\hyperref[tab:8.5]{表 8.5} 中人口普查数据作因子分析. 从 $\mathbf{R}$ 着手, 求极大似然和主成分解. 评论一下你选择的 $m$ 值. 你的分析需包含因子旋转和因子得分的计算.

\noindent 9.25\quad 做\hyperref[tab:4.3]{表 4.3} 给出的及例 4.14 讨论过的“硬度”测量值的因子分析, 计算因子得分并查找数据的离群值. 利用样本协方差矩阵 $\mathbf{S}$.

\noindent 9.26\quad 见例 8.6 的老鼠重量数据, 由样本协方差矩阵着手(见练习 8.15 的 $\sqrt{s_{ii}}$)
(a) 求 $m=1$ 和 $m=2$ 的因子模型的主成分解.
(b) 求 $m=1$ 和 $m=2$ 时载荷和特殊方差的极大似然估计.
(c) 把(a)和(b)的解作最大方差旋转.
